%=============================================================================!
% 1.2  VELOCITY                                                               !
%=============================================================================!
\section{Velocity}
\label{sec:mo-velocity}

Consider a ball thrown vertically upwards. We measure its height $x$ above
the hand, taking upwards as positive, and start the clock at the instant
$t = 0$ at which the ball leaves the hand. Provided that air resistance is
negligible, the heights recorded on a film of the throw are described by
\begin{equation}
   x(t) = v_0 t - \tfrac12 g t^2,
   \label{eq:mo-ball}
\end{equation}
where $v_0$ is the speed at which the ball is thrown and $g$ is a constant,
the same for every object near the Earth's surface. The value of $g$
varies slightly from place to place, and we use its conventional standard
value, \SI{9.80665}{\metre\per\second\squared}, read `metres per second
squared': a metre divided by a second twice, so that $g$ has dimension
$\mathsf{L}\mathsf{T}^{-2}$ (\cref{sec:mo-acceleration} explains why
such a unit arises). For the present,
\cref{eq:mo-ball} is simply a description of what is observed, and
Chapter~2 derives it. Its two terms balance in dimension, as they must:
$v_0 t$ has dimension $\mathsf{L}\mathsf{T}^{-1}\times\mathsf{T} =
\mathsf{L}$, and $g t^2$ has dimension $\mathsf{L}\mathsf{T}^{-2}\times
\mathsf{T}^2 = \mathsf{L}$. The figures and worked numbers of this
chapter use $v_0 = \SI{12}{\metre\per\second}$.

\subsection*{Average velocity}

To find how fast the ball is moving at a particular time $t_1$, chosen
anywhere during the flight, we begin with a finite interval, from $t_1$ to
a slightly later time $t_1 + \tau$. The instant $t_1$ is not the moment of
the throw, which is $t = 0$; throughout the book a subscript 0 marks a
value at that starting moment, as in $v_0$, and other subscripts mark
instants chosen later.
During this interval the ball rises by $x(t_1+\tau) - x(t_1)$, and
dividing the rise by the duration $\tau$ gives the average velocity over
the interval. On the graph of $x(t)$ the average velocity is the slope of
the straight line, or secant, joining the points at $t_1$ and $t_1+\tau$
(\cref{fig:mo-secant}a).

For the ball we can compute the average velocity exactly. The height at
the later time follows from \cref{eq:mo-ball} by replacing $t$ with
$t_1+\tau$, and the square is then multiplied out:
\begin{align*}
   x(t_1+\tau) &= v_0\,(t_1+\tau) - \tfrac12 g\,(t_1+\tau)^2
      && \text{replacing $t$ by $t_1+\tau$}\\
   &= v_0 t_1 + v_0\tau - \tfrac12 g\,\bigl(t_1^2 + 2t_1\tau + \tau^2\bigr)
      && \text{expanding the square}\\
   &= v_0 t_1 - \tfrac12 g t_1^2 + v_0\tau - g t_1\tau - \tfrac12 g\tau^2
      && \text{multiplying out and collecting terms.}
\end{align*}
The first two terms of the last line are exactly $x(t_1)$, so subtracting
$x(t_1)$ leaves only the terms that contain $\tau$,
\begin{equation*}
   x(t_1+\tau) - x(t_1) = v_0\tau - g t_1\tau - \tfrac12 g\tau^2 ,
\end{equation*}
and dividing by $\tau$ gives the average velocity,
\begin{equation}
   \frac{x(t_1+\tau) - x(t_1)}{\tau} = v_0 - g t_1 - \tfrac12 g\tau .
   \label{eq:mo-secant}
\end{equation}
The average velocity depends on the length of the interval through its
final term, $-\tfrac12 g\tau$; the first two terms are fixed once $t_1$ is
chosen. With $v_0 = \SI{12}{\metre\per\second}$ and $t_1 =
\SI{0.4}{\second}$, for example, when the ball is \SI{4.015}{\metre} up
and still rising, substituting into \cref{eq:mo-secant}
gives, in metres per second with $\tau$ in seconds,
\begin{equation*}
   v_0 - g t_1 - \tfrac12 g\tau
   = 12 - 9.80665\times0.4 - \tfrac12\times9.80665\times\tau
   = 8.077 - 4.903\,\tau .
\end{equation*}
Each interval therefore lowers the average below \SI{8.077}{\metre\per
\second} by $4.903\,\tau$, and the table lists the result for three
intervals, each worked out before rounding:
\begin{center}
\begin{tabular}{ccc}
\toprule
$\tau$ / \si{\second} & $\tfrac12 g\tau$ / \si{\metre\per\second}
   & average velocity / \si{\metre\per\second}\\
\midrule
1.2 & 5.884 & 2.193\\
0.6 & 2.942 & 5.135\\
0.2 & 0.981 & 7.097\\
\bottomrule
\end{tabular}
\end{center}
The shorter the interval, the closer the average comes to
\SI{8.077}{\metre\per\second} (\cref{fig:mo-secant}b). The reason is that
the velocity of the ball, counted positive upwards, decreases throughout
each interval: the ball slows as it rises, and in the longest interval,
which runs past the top of the flight at \SI{1.22}{\second} (a time
found in \cref{sec:mo-acceleration}), it then
falls ever faster. Its average velocity is therefore lower than its
velocity at the start, $t_1$; a shorter interval leaves less time for
the velocity to fall, and the average climbs back towards the velocity at
$t_1$ itself. The value $v_0 - g t_1 =
\SI{8.077}{\metre\per\second}$ is therefore the velocity at the instant
$t_1$, as we now make precise.

\begin{figure}[tb]
\centering
\includegraphics{ch01_motion/secant}
\caption{Velocity as a limit. (a) Height of the ball of
\cref{eq:mo-ball} against time, with secants from $t_1 = \SI{0.4}{\second}$
over three intervals $\tau$ and the tangent at $t_1$. (b) The slope of the
secant against $\tau$, which falls short of $v(t_1)$ by $g\tau/2$ and
reaches it as $\tau \to 0$.}
\label{fig:mo-secant}
\end{figure}

\subsection*{Velocity at an instant}

The velocity at the single instant $t_1$ cannot be found by setting $\tau
= 0$ in the quotient on the left of \cref{eq:mo-secant}, because both its
numerator and its denominator would then vanish and $0/0$ has no value.
Instead we ask what value the average approaches as the interval is made
shorter and shorter. In \cref{eq:mo-secant} the answer is plain: the term
$-\tfrac12 g\tau$ can be made as small as we please by choosing $\tau$
small enough, while the first two terms do not depend on $\tau$ at all,
so the average approaches $v_0 - g t_1$. This value is the limit of the
average as $\tau$ tends to zero, and we define it to be the velocity at
$t_1$,
\begin{equation*}
   v(t_1) = \lim_{\tau\to0}\frac{x(t_1+\tau) - x(t_1)}{\tau}
   = \frac{\dd x}{\dd t}\bigg|_{t_1} ,
\end{equation*}
where the vertical bar with $t_1$ means that the derivative is evaluated
at the instant $t = t_1$.
The velocity is therefore the derivative of the position with respect to
time. Geometrically, as $\tau$ shrinks the secant turns about the point
at $t_1$ until it lies along the tangent there (\cref{fig:mo-secant}a), so
the velocity is the slope of the tangent to the graph of $x(t)$. For the
ball, since $t_1$ can be any time,
\begin{equation*}
   v(t) = v_0 - g t ,
\end{equation*}
and at $t_1 = \SI{0.4}{\second}$ the ball is moving upwards at
\SI{8.077}{\metre\per\second}.

\begin{toolbox}[Derivatives of powers]
The calculation above is an instance of a general rule. For a power
$t^n$, with $n$ a whole number $1, 2, 3, \dots$, multiply out $(t+\tau)^n$, the product of $n$ identical factors
$(t+\tau)$. Each term of the result takes either $t$ or $\tau$ from every
factor. Taking $t$ from all $n$ factors gives $t^n$; taking $\tau$ from
exactly one factor, which can be any of the $n$, gives $n$ terms each
equal to $t^{n-1}\tau$; every other choice takes $\tau$ at least twice.
Hence
\begin{equation*}
   (t+\tau)^n = t^n + n t^{n-1}\tau
   + \tau^2\times(\text{terms in } t \text{ and } \tau) .
\end{equation*}
Subtracting $t^n$, dividing by $\tau$ and letting $\tau\to0$ removes
every term that still contains $\tau$, leaving
\begin{equation*}
   \frac{\dd}{\dd t}\,t^n = n\,t^{n-1}.
\end{equation*}
The derivative of a sum is the sum of the derivatives, and a constant
factor passes through unchanged, because the difference quotient has both
properties for every $\tau$. A constant $c$ has derivative zero, since its
difference quotient is $(c - c)/\tau = 0$ for every $\tau$. Applied to \cref{eq:mo-ball}, these rules
give $\dd(v_0 t)/\dd t = v_0$ and $\dd(\tfrac12 g t^2)/\dd t = \tfrac12 g
\times 2t = gt$, and hence $v = v_0 - gt$ once more.
\end{toolbox}

The sign of the velocity carries the direction of motion: it is positive
while the ball rises and negative while it falls. Its magnitude $\lvert
v\rvert$ is the speed, which says how fast the ball moves without saying
in which direction.

\notebook{01}{shrink the interval and watch the secant turn into the tangent}
